% 来源：OEBPS/Text/01_48.xhtml（OPF spine 顺序；仅第一册）
\cleardoublepage
\BookEntry[约翰·班扬《天路历程》]{约翰·班扬《天路历程》}

\BookCover{images/书单32天路历程.jpg}{[英] 约翰·班扬著| 赵沛林、陈亚珂译| 吉林出版集团2011年版}

\begin{BookSummaryBox}
{\centering \textbf{17世纪英国牧师、散文家\\ 约翰·班扬的寓言式小说}\par}
\end{BookSummaryBox}

\subsubsection*{这本书谈什么？}

《天路历程》一书，借用了寓言和梦境的形式，叙述者在梦中看到一个叫作“基督徒”的人正在读一本书，知道了自己居住的城市将遭天火焚毁，惊恐不已。这时“传道者”指点他必须逃离自己的故乡，前往天国才可免于灾祸。基督徒背负着世界的重担，从此踏上了艰难而勇敢的历程，为自己、也为他人寻找救赎。
\par

本书分上、下两卷，最早于1684年出版，写的是基督徒的妻子“女基督徒”和孩子们在一个叫作“无畏”的人的指引下，前往天堂的朝圣过程。
\par

\subsubsection*{芒格为什么谈起这本书？}

芒格在南加州大学GOULD法学院毕业典礼的演讲中，谈了许多如何工作、如何生活的方法。在最后，他说希望这些老人的废话对毕业生们来说是有用的。最后，他用《天路历程》中那位真理剑客年老之后唯一可能说出的话来结束这次演讲：
\par

\begin{quotation}
我的剑传给能挥舞它的人。
\par
\end{quotation}

1678年出版的《天路历程》是约翰·班扬（John Bunyan）所著的寓言故事，它被认为是英语文学中最重要的作品之一。《天路历程》记录的是叙事者在梦中跟随一个基督徒，走过许多虚构的地方：万念俱灰之沼泽、美不胜收之宫殿、卑贱低下之山谷、死神阴影之山谷、虚荣浮夸之市集和犹疑猜忌之堡垒，最终到达他要寻找的天空之城。真理剑客是一位武功高强的朝圣者，独自打败了三个匪徒，在那位基督徒快要到达终点的时候加入了他的朝圣之路。
\par

\subsubsection*{《天路历程》作者约翰·班扬如是说}

“有一个永存的国度可以居住，永恒的生命可以赐给我们，使我们永久居住在这个国度里。”
\par

“有荣耀的冠冕可以给我们，还有会使我们像天空中的太阳那样发亮的衣服。”
\par

“那儿没有哭泣，也没有悲伤，因为那地方的主人会拭掉我们所有的眼泪。”
\par

“我们会跟天使们作伴，看着他们就使你眼花缭乱。你还会碰到成千上万在我们之前到那地方去的人；他们都不会伤害人，而是仁爱的、圣洁的；每个人都在上帝面前生活着，永远得到他的欢心。总之，在那儿我们会看到戴着金冠的长老，携着金竖琴的圣处女，和那些为了对那儿的主人的爱而被世人千刀万剐、烧死、给野兽吃掉，以及淹死在海里的人。他们都很好，好像穿上了永生的衣服。”
\par
